% generated from tools/ladder_report.py CRITERIA: the same table as in every cb ladder report
\section{评测维度与打分标准}
各维度分开报告，不合成总分；表格中 ↑ 越大越好，↓ 越小越好，没有箭头的只作描述。
{\small\begin{longtable}{@{}p{1.6cm}p{2.6cm}p{10.3cm}p{0.9cm}@{}}
\toprule 维度 & 指标 & 计算方法 & 方向 \\ \midrule \endhead
正确性 & 测试点得分 $s_t$ & 每个测试点 0–1：USACO 忽略空白逐 token 比对，相同为 1、否则 0；EGOI、IOI 取官方 checker 给出的 0–1 分；FrontierCS 取官方 checker 输出的 Ratio（达到参考解为 0.6，超过参考解按对数加分，最高 1）。程序编译失败、运行出错、超时或超内存，该测试点为 0。 &  \\
正确性 & 题目得分 $S$（0–100） & USACO：$S = 100 \cdot \frac{1}{T}\sum_t s_t$（每个测试点等分）；EGOI、IOI：$S = 100 \cdot \frac{\sum_g w_g \min_{t \in g} s_t}{\sum_g w_g}$（子任务 $g$ 取其中测试点的最低分，再乘以子任务分值 $w_g$）；FrontierCS：$S = 100 \cdot \frac{1}{T}\sum_t \mathrm{Ratio}_t$。 & ↑ \\
正确性 & 解出 & 所有测试点 $s_t = 1$。 &  \\
正确性 & pass@1、平均得分、pass@$n$ & pass@1 = 解出的采样数 ÷ 采样总数 × 100；平均得分 = 所有采样的 $S$ 的平均；pass@$n$ = 至少一次采样解出的题数 ÷ 题数 × 100。每档采样次数：L1、L2、L5 各 1 次，L3、L4 各 3 次，temperature 0.2。 & ↑ \\
正确性 / 效率 & 宽时限与官方时限 & 每个测试点只运行一次，时限为官方时限的 5 倍（宽时限）。宽时限口径：按上面的规则计分。官方时限口径：CPU 时间超过官方时限的测试点记 0，再按同样的规则算 $S$、pass@1。前者回答“算法对不对”，后者回答“是否和人类选手一样在时限内通过”。 & ↑ \\
效率 & 运行时间倍数 & 只对解出的答案、且官方解在某个测试点上至少需要 0.05 秒 CPU 的题：取官方解最慢的 3 个测试点，答案与官方解交替各运行 3 次取中位数，$r$ = 答案 CPU 时间 ÷ 官方解 CPU 时间（取官方解最慢的那个测试点）；各档报告 $r$ 的几何平均和题数。同语言比较。 & ↓ \\
效率 & 内存倍数 & 解出的答案在所有测试点上的峰值内存 ÷ 官方解的峰值内存，各档取几何平均。 & ↓ \\
复杂度 & 增长幂次 $k$ 与 $\Delta k$ & 以测试点输入字节数为规模 $x$，CPU 时间（或峰值内存）扣除空程序开销为 $y$，在 $y$ 高于噪声下限（10 ms / 2 MB）的测试点上拟合 $\log y = a + k\log x$；在答案与官方解都可测的测试点上拟合 $\log(y_{\text{答案}}/y_{\text{官方}})$ 对 $\log x$ 的斜率 $\Delta k$。$\Delta k \le -0.3$ 为更优（增长慢于官方解，列出供人工复核），$|\Delta k| < 0.3$ 为相同，$\Delta k \ge 0.3$ 为更差；输入大小变化不足 8 倍或可测测试点少于 4 个记为不可测。时间、空间分别判定。 &  \\
失败类型 & 每份未解出答案的主要原因 & 按顺序取第一个成立的：无代码、编译错误、答案错误、运行错误、超内存、超时、部分分（部分测试点得分）；宽时限下解出但官方时限下没有全部通过的记为“只在官方时限下超时”。 &  \\
静态质量 & Pylint、Radon、SonarQube（只限 L1–L2 的 Python 答案） & 按 arXiv:2610.19167 的口径与版本：Pylint 3.3.5 默认配置的问题数（C/R/W/E）与 0–10 分，Radon 6.0.1 的圈复杂度（所有条目之和）、可维护性指数 MI、SLOC，SonarQube 10.4 的可维护性/可靠性/安全问题数；报告正确答案的每份平均与每百行 SLOC 的问题数。 & 见表头 \\
成本与时延 & token、费用、延迟、首个 token 时间 & 费用 = (输入 − 缓存命中) × 输入单价 + 缓存命中 × 命中单价 + 输出（含思考）× 输出单价，单价按每百万 token，分时计价按每次调用的开始时间；延迟为每次调用从发出到收完的时间。 & ↓ \\
泄露检查 & 2026 子集与 USACO 子集 & L3、L4 中 EGOI/IOI 2026 新题与同档 USACO 旧题分开报告 pass@1 与平均得分；新题明显更差时提示旧题可能见过。只作提示，不作结论。 &  \\
\bottomrule
\end{longtable}}
